\documentclass[11pt]{amsart}

\usepackage[T1]{fontenc}
\usepackage{lmodern}
\usepackage{microtype}
\usepackage{mathtools,amssymb,amsthm}
\usepackage{array}
\usepackage[hidelinks]{hyperref}

\hypersetup{
  pdftitle={A Hash-and-Separable Matrix of Order 9}
}

\newtheorem{theorem}{Theorem}
\newtheorem{lemma}[theorem]{Lemma}
\newtheorem{corollary}[theorem]{Corollary}
\theoremstyle{definition}
\newtheorem*{definition}{Definition}
\theoremstyle{remark}
\newtheorem*{remark}{Remark}

\DeclareMathOperator{\col}{col}
\DeclareMathOperator{\row}{row}
\newcommand{\sh}{\operatorname{sh}}
\newcommand{\st}{\ast}

\title[A hash-and-separable matrix of order 9]
{A Hash-and-Separable Matrix of Order~9}

\date{}

\subjclass[2020]{94B25, 05B30, 05D05}
\keywords{hash code, separable code, superimposed code, screening design,
characteristic matrix, explicit construction}

\begin{document}

\begin{abstract}
D'yachkov's lecture notes on designing screening experiments print a
$C_{HS}(13,4)$-matrix and then ask, as Open problem~1, whether a
$C_{HS}(q,4)$-matrix can be constructed with a smaller $q$. We exhibit a
$C_{HS}(9,4)$-matrix, so that $q_4^{HS}\le 9$. The matrix is
printed in full; its three defining conditions reduce to nine $4\times4$ blocks
and to $162$ configurations on its $84$ triples of rows. We claim no priority:
whether such a matrix is already recorded in the literature is not determined
here, and neither is whether $q_4^{HS}=9$, which would need non-existence at
$q=7$ and $q=8$.
\end{abstract}

\maketitle

\section{The question}

Throughout, $q\ge k\ge 2$ are integers, $[q]=\{1,\dots,q\}$, and $\st$ is a
symbol not in $[q]$.

\begin{definition}
A \emph{$C(q,k)$-matrix} is a $q\times q$ matrix $C=(c_i(j))_{i,j\in[q]}$ with
entries in $\{\st\}\cup[q]$.
\end{definition}

Such a matrix is the characteristic matrix of a code of length $3$: its
non-$\st$ cells are read as the codewords $(c_i(j),i,j)$, and D'yachkov's
Propositions~5, 6 and~7 \cite{Dya} translate the three properties that code may
have into conditions on the grid. We reproduce those three conditions in the
form in which they are used below.

\begin{enumerate}
\item[(P5)] Every symbol $a\in[q]$ occurs exactly $k$ times in $C$; every row
and every column of $C$ contains exactly $q-k$ stars; and no symbol occurs
twice in one row or twice in one column.

\item[(P6)] If $c_i(j)=c_m(n)=a\ne\st$ with $i\ne m$ and $j\ne n$, then
$c_i(n)=c_m(j)=\st$.

\item[(P7)] No $3\times3$ submatrix of $C$, on rows $i_1<i_2<i_3$ and columns
$j_1<j_2<j_3$, has any of the six forms obtained by permuting the columns of
\[
\begin{pmatrix}
\st & a & c\\
a & \st & b\\
c & b & \st
\end{pmatrix},
\qquad a,b,c\in[q]\ \text{pairwise distinct.}
\]
\end{enumerate}

By Proposition~6, a $C(q,k)$-matrix satisfying (P5) and (P6) is exactly a
\emph{$C_H(q,k)$-matrix} (the hash case); by Proposition~7, one satisfying (P5),
(P6) and (P7) is exactly a \emph{$C_{HS}(q,k)$-matrix}. The source supplies no
grid condition for separable codes alone, and nothing below needs one:
separability enters only through Definition~4 of the source, applied directly to
a code, in the remark of Section~2. Following \cite{Dya} we write
$q_k^{HS}$ for the least $q\ge k$ admitting a $C_{HS}(q,k)$-matrix, that text
asking verbatim
\begin{quote}
``How to find the minimal possible integer $q_k\ge k$ such that there exists
$C_S(q_k,k)$-matrix, $C_H(q_k,k)$-matrix or $C_{HS}(q_k,k)$-matrix?''
\end{quote}
Example~10 of \cite{Dya} then reads ``Let $k=4$. For hash codes, $q_4=8$ and
for hash\&separable codes, $q_4=13$'', and prints a $C_H(8,4)$ and a
$C_{HS}(13,4)$ to support the two values. Immediately afterwards comes the
display containing the question at issue:
\begin{quote}
``\textbf{Open problems}. \textbf{1}. Is it possible to construct a
$C_{HS}(q,4)$-matrix if $q<13$? \textbf{2}. Is it possible to construct
$C_{HS}(q,k)$-matrices, if $k\ge5$ and $q<k^2$?''
\end{quote}

\medskip
\noindent\textbf{Source.}
Quotations are from the file \texttt{lect\_dse.tex} in the e-print source of
\texttt{arXiv:1401.7505v1}; the printed lecture notes may number the
propositions, examples and open problems differently, and nothing below depends
on the labelling, since (P5)--(P7) and the question are reproduced in full
above. The value $13$ is the smallest order at which the source prints a
$C_{HS}(q,4)$-matrix, and Open problem~1 asks, one paragraph later, for a
smaller one; the source does not prove that none exists.

\medskip
\noindent\textbf{Related objects.}
Conditions (P5) and (P6) by themselves describe an object studied elsewhere
under another name: a $q\times q$ array in which every row, every column and
every symbol carries exactly $k$ entries, and in which two equal entries in
distinct rows and distinct columns force the two opposite cells to be empty,
matches the definition of a \emph{Blackburn $k$-plex} in \cite{Sto}, which
proves a lower bound on the order of such an array and describes constructions
for $k\le5$; (P7) is not among the conditions treated there. F\"uredi and
Ruszink\'o \cite{FR} characterise $3$-union-free $3$-partite linear hypergraphs
by the avoidance of three grids, a counterpart of (P7); they give no value of
$q_4^{HS}$, and they name small-value tables of Kim--Lebedev and of Lebedev for
this family. The source states that the results of the section quoted above were
to appear in \cite{DyaRykov}. We claim no priority: a $C_{HS}(q,4)$-matrix with
$q<13$ may already be recorded in that literature. What is offered here is the
matrix and its verification.

\section{The matrix}

\begin{theorem}\label{thm:main}
The $9\times9$ matrix
\[
M=
\begin{pmatrix}
1&2&3&4&\st&\st&\st&\st&\st\\
5&\st&\st&\st&9&\st&4&\st&8\\
6&\st&\st&\st&\st&8&3&2&\st\\
7&\st&9&\st&\st&4&\st&\st&2\\
\st&8&6&\st&1&\st&\st&4&\st\\
\st&\st&\st&9&\st&1&\st&7&3\\
\st&9&\st&\st&\st&\st&1&5&6\\
\st&7&\st&6&3&5&\st&\st&\st\\
\st&\st&5&8&2&\st&7&\st&\st
\end{pmatrix}
\]
is a $C_{HS}(9,4)$-matrix. Consequently $q_4^{HS}\le 9$: a $C_{HS}(q,4)$-matrix
with $q<13$ exists, which is what Open problem~1 of \cite{Dya} asks for.
\end{theorem}

For $a\in[q]$ let $R_a$ and $C_a$ be the sets of rows and of columns of $M$
carrying $a$, and, for $i\in R_a$, let $\col_a(i)$ be the column of $a$ in
row~$i$. Reading $M$ off:
\[
\renewcommand{\arraystretch}{1.05}
\begin{array}{r|l|l}
a & R_a & C_a\\\hline
1 & 1,5,6,7 & 1,5,6,7\\
2 & 1,3,4,9 & 2,5,8,9\\
3 & 1,3,6,8 & 3,5,7,9\\
4 & 1,2,4,5 & 4,6,7,8\\
5 & 2,7,8,9 & 1,3,6,8\\
6 & 3,5,7,8 & 1,3,4,9\\
7 & 4,6,8,9 & 1,2,7,8\\
8 & 2,3,5,9 & 2,4,6,9\\
9 & 2,4,6,7 & 2,3,4,5
\end{array}
\]

\begin{proof}[(P5)]
Each row and each column of $M$ carries exactly $4$ symbols, hence exactly
$9-4=5$ stars, and the four symbols in a row, and in a column, are distinct.
Each $a\in[9]$ occupies the four cells listed in the table, so each occurs
exactly $4$ times; the $36=4\cdot9$ non-star cells are accounted for exactly
once.
\end{proof}

The next lemma is what makes (P6) a hand check: it replaces the quantifier over
quadruples $(i,j,m,n)$ by one $k\times k$ block per symbol.

\begin{lemma}\label{lem:block}
Let $C$ be a $C(q,k)$-matrix satisfying \textup{(P5)} and let $a\in[q]$. Then
\textup{(P6)} holds for the symbol $a$ if and only if the $k\times k$ submatrix
$C[R_a\times C_a]$ contains no non-star entry other than the $k$ cells of $a$
itself, i.e.\ if and only if that submatrix is a $k\times k$ permutation matrix
over $\{a,\st\}$.
\end{lemma}

\begin{proof}
By (P5) the $k$ cells of $a$ lie in $k$ distinct rows and $k$ distinct columns,
so $i\mapsto\col_a(i)$ is a bijection $R_a\to C_a$; write $\row_a$ for its
inverse. Let $i\in R_a$ and $n\in C_a$ with $n\ne\col_a(i)$, and put
$m=\row_a(n)$. Then $m\ne i$, and $c_i(\col_a(i))=c_m(n)=a$ with
$\col_a(i)\ne n$, so (P6) forces $c_i(n)=\st$. Every cell of
$C[R_a\times C_a]$ off the bijection is of this form, which gives one
direction. Conversely, any instance of (P6) for the symbol $a$ has $j=\col_a(i)$
and $n=\col_a(m)$ with $i,m\in R_a$ and $j,n\in C_a$, so the two cells
$c_i(n)$ and $c_m(j)$ it constrains lie in $C[R_a\times C_a]$ off the
bijection.
\end{proof}

\begin{proof}[(P6)]
By Lemma~\ref{lem:block} it suffices to inspect nine $4\times4$ blocks. For
$a=1$, $R_1=C_1=\{1,5,6,7\}$ and
$M[\{1,5,6,7\}\times\{1,5,6,7\}]$ has the four $1$'s on its diagonal and stars
elsewhere. For $a=2$, $R_2=\{1,3,4,9\}$, $C_2=\{2,5,8,9\}$, and the block
carries $2$ at $(1,2),(3,8),(4,9),(9,5)$ and stars in its other twelve cells.
The remaining seven blocks are read off the same way from the table; each is a
permutation matrix.
\end{proof}

For (P7) we first remove the star conditions, which (P6) already implies. Write
$\sh(i,m)$ for the set of symbols occurring in both row $i$ and row $m$.

\begin{lemma}\label{lem:red}
Let $C$ be a $C(q,k)$-matrix satisfying \textup{(P5)} and \textup{(P6)}. Then
$C$ fails \textup{(P7)} if and only if there are rows $i_1<i_2<i_3$ and
pairwise distinct symbols
\[
a\in\sh(i_1,i_2),\qquad b\in\sh(i_2,i_3),\qquad c\in\sh(i_1,i_3)
\]
with
\begin{equation}\label{eq:cols}
\col_a(i_1)=\col_b(i_3),\qquad
\col_a(i_2)=\col_c(i_3),\qquad
\col_b(i_2)=\col_c(i_1),
\end{equation}
these three columns being distinct.
\end{lemma}

\begin{proof}
Since the six forms of (P7) are the column permutations of one, and that one is
symmetric, an occurrence on a set of three rows and a set of three columns may
be taken with the rows in increasing order and the columns free. Such an
occurrence, with columns $j_1,j_2,j_3$, has $a$ at $(i_1,j_2)$ and $(i_2,j_1)$,
$b$ at $(i_2,j_3)$ and $(i_3,j_2)$, and $c$ at $(i_1,j_3)$ and $(i_3,j_1)$;
reading off the columns gives \eqref{eq:cols} with
$(j_1,j_2,j_3)=(\col_a(i_2),\col_a(i_1),\col_b(i_2))$, distinct. Conversely,
suppose \eqref{eq:cols} holds with $j_1,j_2,j_3$ as just named. The six named
cells are then filled as required, and the three diagonal cells are stars: $a$
sits at $(i_1,j_2)$ and $(i_2,j_1)$ with $i_1\ne i_2$ and $j_1\ne j_2$, so (P6)
gives $c_{i_1}(j_1)=c_{i_2}(j_2)=\st$; applying (P6) to $b$ gives
$c_{i_2}(j_2)=c_{i_3}(j_3)=\st$ and to $c$ gives
$c_{i_1}(j_1)=c_{i_3}(j_3)=\st$. So the submatrix on those rows and columns is
a forbidden form.
\end{proof}

\begin{proof}[(P7)]
$M$ has $\binom92=36$ pairs of rows; exactly $18$ of them share two symbols and
$18$ share exactly one, a total of $2\cdot18+1\cdot18=54=q\binom k2$ shared
incidences. In particular every one of the $\binom93=84$ triples of rows is
pairwise sharing, so Lemma~\ref{lem:red} must be applied to all of them. Running
over the $84$ triples and over the pairwise distinct choices of
$(a,b,c)\in\sh(i_1,i_2)\times\sh(i_2,i_3)\times\sh(i_1,i_3)$ leaves exactly
$162$ candidate configurations, and each of the $162$ violates one of the three
equations \eqref{eq:cols}. Hence (P7) holds.
\end{proof}

The $162$ configurations are enumerated by the program described in
Section~\ref{sec:repro}, which also verifies (P7) on $M$ with no reduction at
all and re-derives every count quoted above. This completes the proof of
Theorem~\ref{thm:main}.

\begin{remark}[the definitions applied directly]
Read $M$'s $36$ non-star cells as the codewords $(c_i(j),i,j)$. The resulting
code of length $3$ satisfies Definitions~2, 4 and~5 of \cite{Dya} directly:
each of the nine symbols occurs exactly $4$ times in each of the three
coordinates, any two codewords agree in at most one coordinate, every one of the
$\binom{36}3=7140$ triples has a coordinate on which all three differ, and the
coordinatewise unions of the $7806$ subsets of size at most $3$ are pairwise
distinct. So $M$ is a $C_{HS}(9,4)$-matrix by the raw definitions and not only
by way of Propositions~5--7.
\end{remark}

\section{What is and is not settled}

\begin{lemma}\label{lem:2k}
Every $C_H(q,k)$-matrix, hence every $C_{HS}(q,k)$-matrix, has $q\ge2k-1$.
\end{lemma}

\begin{proof}
Fix $a\in[q]$. By Lemma~\ref{lem:block} each of the $k$ rows in $R_a$ already
carries $k-1$ stars inside $C[R_a\times C_a]$, while by (P5) a row carries
exactly $q-k$ stars in total. So $k-1\le q-k$.
\end{proof}

For $k=4$ Lemma~\ref{lem:2k} gives $q\ge7$, so $q=7$ and $q=8$ are the orders
below $9$ that this note does not decide.

What is settled is $q_4^{HS}\le9$, by finite checks on the printed matrix. The
exact value is not: $q_4^{HS}=9$ would additionally require that no
$C_{HS}(7,4)$- and no $C_{HS}(8,4)$-matrix exists, and neither non-existence is
attempted here or by the accompanying program. Open problem~2 of \cite{Dya}
($k\ge5$, $q<k^2$) is untouched.

\section{Reproduction and verification}\label{sec:repro}

Nothing in Theorem~\ref{thm:main} needs a computer to be \emph{stated}: the
matrix is printed above, (P5) and (P6) are read off it by way of the table and
Lemma~\ref{lem:block}, and (P7) reduces by Lemma~\ref{lem:red} to $162$
configurations. The program \texttt{verify.py} in this folder re-derives it all
from the printed object, run as \texttt{python3\ verify.py} with the standard
library only and in
exact integer arithmetic: (P5) clause by clause; (P6) both by exhaustive search
over the $2304$ quadruples $(i,j,m,n)$ and by the nine blocks of
Lemma~\ref{lem:block}; (P7) both by the reduction of Lemma~\ref{lem:red} ($162$
candidates, none surviving) and by unreduced search over row and column triples;
and the definitional route of the remark above, including all $7140$ triples and
all $7806$ subsets. Each check is also run on the labelled matrices of
Examples~7--10 of \cite{Dya}, which the source itself marks as satisfying some
of the three conditions and failing others. The treatments of (P7) are encodings
inside one program and share one transcription of $M$, so their agreement bears
on the encodings and is not an independent confirmation. The recorded run
\texttt{verify.output.txt} reports $50$ checks, all passing, and lists in its
own output what it does not cover --- in particular the non-existence claims at
$q=7$ and $q=8$, which it does not attempt.

\begin{thebibliography}{9}

\bibitem{Dya}
A.~G. D'yachkov,
\emph{Lectures on Designing Screening Experiments},
POSTECH Combinatorial and Computational Mathematics Center Lecture Note
Series \textbf{10} (February 2004);
\href{https://arxiv.org/abs/1401.7505v1}{arXiv:1401.7505v1} (29 January 2014).
All quotations in this note are from \texttt{lect\_dse.tex} in the v1 e-print
source.

\bibitem{DyaRykov}
A.~G. D'yachkov and V.~V. Rykov,
\emph{Optimal superimposed codes and designs for R\'enyi's search model},
J. Statist. Plann. Inference \textbf{100} (2001), no.~2, 281--302.
\href{https://doi.org/10.1016/S0378-3758(01)00140-9}{doi:10.1016/S0378-3758(01)00140-9}.

\bibitem{FR}
Z.~F\"uredi and M.~Ruszink\'o,
\emph{Uniform hypergraphs containing no grids},
Adv. Math. (2013);
\href{https://arxiv.org/abs/1103.1691}{arXiv:1103.1691}.

\bibitem{Sto}
R.~J. Stones,
\emph{$k$-plex $2$-erasure codes and Blackburn partial Latin squares},
IEEE Trans. Inform. Theory \textbf{66} (2020), no.~6, 3704--3713.
\href{https://doi.org/10.1109/TIT.2020.2967758}{doi:10.1109/TIT.2020.2967758}.

\end{thebibliography}

\end{document}
